\documentclass[11pt,letterpaper]{article}
\usepackage[T1]{fontenc}
\usepackage{lmodern}
\usepackage{amsmath,amssymb,amsthm}
\usepackage[margin=0.91in,headheight=14pt]{geometry}
\usepackage{microtype}
\usepackage{enumitem}
\usepackage{booktabs}
\usepackage{fancyhdr}
\usepackage{xcolor}
\usepackage{hyperref}
\hypersetup{colorlinks=true,linkcolor=black,citecolor=black,urlcolor=blue!45!black,
pdftitle={Linear Programming by Sparse Certificates},
pdfsubject={A hand calculation approach to exact linear programming certificates},
pdfauthor={Expository companion to Codex for Math}}
\newcommand{\R}{\mathbb R}
\newcommand{\Q}{\mathbb Q}
\newcommand{\T}{\mathsf T}
\DeclareMathOperator{\rank}{rank}
\DeclareMathOperator{\cone}{cone}
\DeclareMathOperator{\im}{im}
\DeclareMathOperator{\supp}{supp}
\DeclareMathOperator{\spanop}{span}
\newtheorem{theorem}{Theorem}[section]
\newtheorem{proposition}[theorem]{Proposition}
\newtheorem{lemma}[theorem]{Lemma}
\newtheorem{corollary}[theorem]{Corollary}
\theoremstyle{definition}
\newtheorem{definition}[theorem]{Definition}
\newtheorem{example}[theorem]{Example}
\theoremstyle{remark}
\newtheorem{remark}[theorem]{Remark}
\numberwithin{equation}{section}
\setlist{itemsep=3pt,topsep=6pt}
\setlength{\parskip}{3pt}
\setlength{\parindent}{1.2em}
\setlength{\emergencystretch}{1.5em}
\allowdisplaybreaks[2]
\pagestyle{fancy}
\fancyhf{}
\fancyhead[L]{\small Linear programming by sparse certificates}
\fancyhead[R]{\small A hand calculation approach}
\fancyfoot[C]{\thepage}
\renewcommand{\headrulewidth}{0pt}
\title{\vspace{-1.0cm}\textbf{Linear Programming by Sparse Certificates}\\[5pt]
\large A Hand Calculation Approach}
\author{An expository companion to the Codex for Math LP studies}
\date{13 October 2026}

\begin{document}
\maketitle
\thispagestyle{plain}
\begin{abstract}
A short nonnegative combination of constraints can prove a linear-programming
bound without introducing slack or artificial decision variables. The main
question is how to find such a combination systematically. We first eliminate
consistent equalities and then work with the projected inequality normals.
If their rank is \(k\), every valid upper bound on a nonempty feasible set has
a certificate using at most \(k\) inequalities. Infeasibility has a different
geometry: a contradiction may require \(k+1\) inequalities. We prove both
statements, explain the necessary hypotheses, and obtain a finite exact
support-enumeration procedure for rational data. Worked calculations show how
to discover the multipliers, recover the original equalities, and distinguish
a bound from a proved optimum. Degeneracy, normalization, nonuniqueness, and
implementation limits receive explicit treatment. The foundations are
classical; the purpose is a transparent method for hand-sized problems,
rather than a new duality theorem or a claim of general solver superiority.
\end{abstract}

\noindent\textbf{Reading guide.}
Sections~\ref{sec:certificates}--\ref{sec:elimination} establish the certificate
language and work a three-variable example. Sections~\ref{sec:existence} and
\ref{sec:sparsity} give the proofs behind sparse discovery.
Sections~\ref{sec:search}--\ref{sec:examples} turn those proofs into a practical
procedure. Section~\ref{sec:evidence} separates mathematical results from the
published computational and formal evidence.

\section{The problem and the meaning of a certificate}
\label{sec:certificates}

Let \(A\in\R^{m\times n}\), \(b\in\R^m\),
\(E\in\R^{p\times n}\), \(f\in\R^p\), and \(c\in\R^n\). We consider
\begin{equation}\label{eq:model}
 P=\{x\in\R^n:Ax\le b,\ Ex=f\},
 \qquad \text{maximize }\ell_c(x)=c^\T x .
\end{equation}
The inequality \(Ax\le b\) is coordinatewise. All variables are free real
variables; a sign restriction such as \(x_j\ge0\) is an explicit row
\(-x_j\le0\). We use only finitely many nonstrict constraints.
Strict inequalities and integer restrictions are outside this model.

The symbol \(\ell_c\) denotes the linear \emph{map}
\(x\mapsto c^\T x\); \(\ell_c(x)\) is a real number.
An optimizer \(x^\star\) is a feasible vector, whereas its objective value
\(v^\star=\ell_c(x^\star)\) is a scalar. A proposed upper bound \(\delta\)
need not be the optimum. Matrix products have their ordinary meanings and
every indexed sum is written explicitly; there is no implicit summation
convention.

A row \(a^\T x\ge b_0\) is normalized by negating \emph{both} sides.
Minimizing \(c^\T x\) is equivalent to maximizing \((-c)^\T x\),
with the final objective value negated. An equality retains an unrestricted
signed multiplier. There is no need to replace it with a pair of inequalities.

\subsection{Weighted addition and weak duality}

\begin{definition}\label{def:bound}
A certificate for the bound \(c^\T x\le\delta\) consists of
\(\lambda\in\R^m\) and \(\mu\in\R^p\) satisfying
\begin{equation}\label{eq:certificate}
 \lambda\ge0,\qquad A^\T\lambda+E^\T\mu=c,\qquad
 b^\T\lambda+f^\T\mu\le\delta.
\end{equation}
The \emph{positive inequality support} is
\(\supp_+(\lambda)=\{i:\lambda_i>0\}\). Its size counts original inequality
rows with nonzero weight, not equality rows, active constraints, or digits.
\end{definition}

\begin{proposition}[Certificate soundness]\label{prop:sound}
If \eqref{eq:certificate} holds, then \(c^\T x\le\delta\) for every
\(x\in P\). If, in addition, \(x^\star\in P\) and
\(c^\T x^\star=\delta\), then \(x^\star\) is an optimizer with value \(\delta\).
\end{proposition}
\begin{proof}
For each feasible \(x\), the coefficient identity gives
\[
 c^\T x=\lambda^\T Ax+\mu^\T Ex
 \le\lambda^\T b+\mu^\T f\le\delta.
\]
The inequality follows from the nonnegativity of each \(\lambda_i\);
the equality terms need no sign restriction because \(Ex=f\).
An attaining feasible point then matches an upper bound valid for all
feasible points.
\end{proof}

This is weak duality in its most direct form. The multiplier pair is feasible
for the dual constraint \(A^\T\lambda+E^\T\mu=c\), \(\lambda\ge0\), and its
dual objective is \(b^\T\lambda+f^\T\mu\). For hand calculation, it is often
more useful to regard these equations as a request: combine a few constraint
left sides to produce the objective \emph{as an identity in all coordinates}.

Matching only the numerical value at one candidate is insufficient. For
example, the constraint \(x\le1\) and objective \(x+y\) agree at \((1,0)\),
but their coefficient vectors differ, and \(y\) may increase without limit.
Likewise, a point attaining a certified bound must still satisfy every row.
Under \(x\le0\) and \(y\ge1\), the point \((0,0)\) attains the bound \(x\le0\)
but is not feasible.

\subsection{The gap and the other two elementary certificates}

For nonnegative inequality multipliers satisfying the coefficient identity,
every \(x\in P\) satisfies
\begin{equation}\label{eq:gap}
 b^\T\lambda+f^\T\mu-c^\T x
   =\sum_{i=1}^m\lambda_i\bigl(b_i-(Ax)_i\bigr).
\end{equation}
All terms on the right are nonnegative. Equality of primal and certificate
values therefore implies
\[
 \lambda_i\bigl(b_i-(Ax)_i\bigr)=0\qquad(1\le i\le m).
\]
This is complementary slackness. A positive weight forces its row to be
active, meaning that its inequality is an equality at that point.
An active row need not have a positive weight.

A \emph{contradiction certificate} satisfies
\begin{equation}\label{eq:contradiction}
 \lambda\ge0,\qquad A^\T\lambda+E^\T\mu=0,\qquad
 b^\T\lambda+f^\T\mu<0.
\end{equation}
Weighted addition would give
\(0\le b^\T\lambda+f^\T\mu<0\), so \(P\) is empty.

For comparison, an \emph{improving-ray certificate} consists of a feasible
base \(x_0\in P\) and a direction \(v\) such that
\begin{equation}\label{eq:ray}
 Av\le0,\qquad Ev=0,\qquad c^\T v>0.
\end{equation}
Then \(x_0+t v\in P\) for every \(t\ge0\), and
\(c^\T(x_0+t v)=c^\T x_0+t\,c^\T v\) tends to \(+\infty\).
The feasible base is essential: homogeneous direction equations alone do not
establish that any feasible point exists. The search developed below concerns
bounds and contradictions, not general discovery of feasible points or rays.

\section{Eliminating equalities before searching}
\label{sec:elimination}

Equalities carry useful information about which directions are actually free.
Using them first can make a three- or ten-variable calculation effectively
one-dimensional. Gaussian elimination is enough; an orthonormal basis is
unnecessary and may be inconvenient for rational arithmetic.

\subsection{Inconsistent equalities must be handled first}

Row reduction of the augmented matrix \([\,E\ f\,]\) decides whether
\(Ex=f\) is consistent. If it is not, the row operations produce a vector
\(h\in\R^p\) with \(E^\T h=0\) and \(f^\T h\ne0\).
Setting
\begin{equation}\label{eq:equality-contradiction}
 \mu=-\frac{h}{f^\T h}
\end{equation}
gives \(E^\T\mu=0\), \(f^\T\mu=-1\): a contradiction using zero
inequalities. For rational data the row operations and \(\mu\) are rational.
This is the finite-dimensional Fredholm alternative; see
Border~\cite[Theorems 3--5, p.~3]{border}.

Do not form an affine parametrization of an inconsistent system. For example,
\(x=0\), \(2x=1\) already yield \(0=-1\) by taking weights \(2,-1\).

\subsection{The reduced model and exact recovery}

Assume henceforth that \(Ex=f\) is consistent. Put \(r=\rank(E)\) and
\(d=n-r\). Choose a particular solution \(x_0\) and a matrix
\(N\in\R^{n\times d}\) whose columns are a basis of \(\ker E\). The affine map
\begin{equation}\label{eq:param}
 \Phi:\R^d\longrightarrow\{x:Ex=f\},\qquad
 \Phi(z)=x_0+Nz
\end{equation}
is bijective. It is a map; \(x_0+Nz\) is its value at \(z\).
Substitution gives
\begin{equation}\label{eq:reduction}
 B=AN,\quad \beta=b-Ax_0,\quad q=N^\T c,\quad
 \eta=\delta-c^\T x_0 .
\end{equation}
Thus \(x=\Phi(z)\) is feasible exactly when \(Bz\le\beta\), and the proposed
bound becomes \(q^\T z\le\eta\). This is the usual equality-elimination
construction~\cite[slide 4.35, PDF p.~122]{boyd}. The dimension here is
\(n-\rank(E)\), allowing dependent equality rows.

Write \(a_i\in\R^d\) for the transpose of row \(i\) of \(B\), and put
\[
 M=B^\T=[\,a_1\ \cdots\ a_m\,],\qquad k=\rank(B).
\]
A reduced bound certificate is simply
\begin{equation}\label{eq:reduced-certificate}
 \lambda\ge0,\qquad M\lambda=q,\qquad \beta^\T\lambda\le\eta.
\end{equation}

\begin{proposition}[Elimination preserves certificates]\label{prop:recovery}
The original bound certificate equations \eqref{eq:certificate} are solvable
with a given \(\lambda\ge0\) if and only if
\eqref{eq:reduced-certificate} holds. Equality multipliers can be recovered
by solving
\begin{equation}\label{eq:recover}
 E^\T\mu=c-A^\T\lambda .
\end{equation}
For every solution of this equation,
\begin{equation}\label{eq:recover-cost}
 b^\T\lambda+f^\T\mu=c^\T x_0+\beta^\T\lambda .
\end{equation}
The analogous statement for contradictions uses the zero objective and a
strictly negative reduced cost.
\end{proposition}
\begin{proof}
Multiplying the original coefficient equation by \(N^\T\) gives
\(M\lambda=q\). Conversely, that equation says
\(N^\T(c-A^\T\lambda)=0\). Since \(\im N=\ker E\), the residual is
orthogonal to \(\ker E\), hence belongs to \(\im E^\T\).
This proves the solvability of \eqref{eq:recover}. Since \(Ex_0=f\),
\[
 f^\T\mu=x_0^\T E^\T\mu
     =c^\T x_0-\lambda^\T Ax_0.
\]
Adding \(b^\T\lambda\) proves \eqref{eq:recover-cost}, and subtracting
\(c^\T x_0\) proves equivalence of the cost tests. For a contradiction the
same calculation uses the zero objective.
\end{proof}

Different recovered vectors \(\mu\) give the same combined constant.
One may solve using a basis of the original equality rows and set the
remaining equality weights to zero. In this consistent-equality recovery,
this uses at most \(r\) equality rows.
Our support bounds count inequalities only.

The relevant rank has a coordinate-free expression:
\begin{equation}\label{eq:intrinsic-rank}
 k=\rank(AN)=\rank\begin{pmatrix}A\\E\end{pmatrix}-\rank(E)\le d.
\end{equation}
Indeed, restrict the covectors in the row space of \(\begin{pmatrix}A\\E\end{pmatrix}\)
to \(\ker E\). The kernel of this restriction is exactly the row space of
\(E\), and its image is the span of the restricted inequality covectors.
Rank-nullity gives \eqref{eq:intrinsic-rank}. Changing the basis \(N\) changes
coordinates and possibly arithmetic size, but not \(k\).

\subsection{A complete hand calculation in three variables}

Maximize \(2x-y+3z\) subject to
\begin{align*}
 x+2y+z&\le10,& -z&\le0,&2x+4y+2z&\le20,\\
 x-y&=1,&y-z&=1,&x-z&=2.
\end{align*}
The last equality is the sum of the first two; the equality rank is \(2\).
Set \(z=t\). Then
\[
 (x,y,z)=(2+t,1+t,t),\qquad
 4t\le6,\quad -t\le0,\quad8t\le12,
\]
and the objective becomes \(3+4t\). The reduced feasible interval is
\(0\le t\le3/2\). Its objective increases with \(t\), so the candidate is
\[
 t^\star=\frac32,\qquad
 (x^\star,y^\star,z^\star)=\left(\frac72,\frac52,\frac32\right),
 \qquad v^\star=9.
\]

Here is how to \emph{discover} the certificate. The projected objective
coefficient is \(4\). The first row has coefficient \(4\), so its multiplier
must be \(1\). Its budget \(6\) gives the reduced bound \(4t\le6\).
No other inequality is needed. Subtract its original normal from the
objective normal:
\[
 (2,-1,3)-(1,2,1)=(1,-3,2)
       =1(1,-1,0)-2(0,1,-1).
\]
This recovers equality weights \((1,-2,0)\). The full identity is
\begin{equation}\label{eq:three-example}
 (x+2y+z)+(x-y)-2(y-z)=2x-y+3z\le10+1-2=9.
\end{equation}
The displayed candidate satisfies all rows and attains \(9\), so the bound is
optimal. The first inequality could instead be replaced by half of the third:
a sparse certificate need not be unique. The reduction explains why the
calculation is short: \(n=3\), but \(d=k=1\).

\section{Why the required certificates exist}
\label{sec:existence}

Weighted addition proves that a certificate is sound. It does not prove that
a certificate exists whenever a bound is valid. To establish that converse,
we use finite conic geometry. The next elementary lemmas also identify the
source of sparsity.

\begin{lemma}[Eliminating conic dependence]\label{lem:conic}
Let \(v_1,\ldots,v_m\) be vectors in a finite-dimensional real vector space.
Every vector \(w=\sum_{i=1}^m\lambda_i v_i\), with \(\lambda_i\ge0\), has
such a representation whose positive-support vectors are linearly independent.
Its support has at most \(\dim\spanop\{v_1,\ldots,v_m\}\) members.
For \(w=0\), the empty representation is permitted.
\end{lemma}
\begin{proof}
If the positive support \(S\) is dependent, choose a nonzero vector \(h\)
supported on \(S\) such that \(\sum_{i\in S}h_i v_i=0\). Reverse its sign
if necessary so that at least one \(h_i\) is positive. Set
\begin{equation}\label{eq:conic-step}
 t=\min_{\substack{i\in S\\h_i>0}}\frac{\lambda_i}{h_i},
 \qquad \lambda'_i=\lambda_i-t h_i.
\end{equation}
For \(h_i>0\), the definition of \(t\) makes \(\lambda'_i\ge0\);
for \(h_i\le0\), it makes \(\lambda'_i\ge\lambda_i\ge0\).
At least one positive entry becomes zero. No new support is introduced
because \(h\) vanishes outside \(S\), and the represented vector is unchanged.
After finitely many steps the support is independent.
\end{proof}

This is the conic form of Carath\'eodory's theorem
\cite[pp.~8--9]{bertsekas3}. Its bound is the dimension of the span,
without the extra one that occurs for convex combinations.

\begin{lemma}[Closedness and separation of a finite cone]\label{lem:separate}
A cone \(C\) generated by finitely many vectors in \(\R^s\) is closed.
If \(y\notin C\), there is \(h\in\R^s\) such that
\[
 h^\T v\ge0\quad(v\in C),\qquad h^\T y<0.
\]
\end{lemma}
\begin{proof}
Let \(y_j\in C\) converge to \(y\). By Lemma~\ref{lem:conic}, each \(y_j\)
has an independent-support representation. There are only finitely many
supports, so a subsequence has one fixed support. If that support is empty,
the subsequence is zero and \(y=0\in C\). Otherwise, the column map on this
support is injective. Its image is a finite-dimensional subspace and thus
closed, and its inverse on that image is continuous. Therefore the
nonnegative coefficient vectors of the subsequence converge to a
nonnegative coefficient vector representing \(y\). This proves closedness.

Choose \(p\in C\) minimizing the Euclidean distance to \(y\). Such a point
exists: a minimizing sequence can be restricted to a bounded closed ball,
and then compactness and closedness supply a limit in \(C\).
For \(v\in C\), the segment from \(p\) to \(v\) lies in \(C\).
Differentiating the squared distance at the segment parameter \(0\) gives
\[
 (y-p)^\T(v-p)\le0 .
\]
Using \(v=0\) and \(v=2p\) gives \((y-p)^\T p=0\).
With \(h=p-y\), it follows that \(h^\T v\ge0\) for all \(v\in C\), while
\[
 h^\T y=-\|y-p\|^2<0.
\]
\end{proof}

\begin{theorem}[The contradiction alternative]\label{thm:farkas}
Exactly one of the following holds:
\begin{enumerate}[label=(\roman*)]
\item there is \(z\in\R^d\) with \(Bz\le\beta\);
\item there is \(\lambda\ge0\) with \(M\lambda=0\) and
      \(\beta^\T\lambda=-1\).
\end{enumerate}
\end{theorem}
\begin{proof}
Both cannot hold, because they would imply
\(0=\lambda^\T Bz\le\lambda^\T\beta=-1\).
Suppose (ii) fails. Then \((0,-1)\) is outside
\(D=\cone\{(a_i,\beta_i):1\le i\le m\}\).
By Lemma~\ref{lem:separate}, there is \((u,t)\in\R^d\times\R\) with
\[
 a_i^\T u+\beta_i t\ge0\quad(1\le i\le m),\qquad -t<0.
\]
Thus \(t>0\), and \(z=-u/t\) satisfies every inequality
\(a_i^\T z\le\beta_i\). This proves (i).
\end{proof}

This is Farkas' alternative for unrestricted variables, in the normalized
form given by Border~\cite[Proposition 35, p.~17]{border}.
A strictly negative certificate can always be rescaled by a positive factor
to make its right side \(-1\).

\begin{theorem}[Existence of bound certificates]\label{thm:existence}
Assume \(Q=\{z:Bz\le\beta\}\ne\varnothing\).
Then \(q^\T z\le\eta\) for every \(z\in Q\) if and only if
\eqref{eq:reduced-certificate} has a solution.
\end{theorem}
\begin{proof}
The reverse implication is weighted addition. For the forward implication,
consider the closed finite cone
\[
 C=\cone\{(a_i,\beta_i):1\le i\le m\}+\cone\{(0,1)\}.
\]
If \((q,\eta)\notin C\), separation gives \((u,t)\) satisfying
\[
 a_i^\T u+\beta_i t\ge0\quad(1\le i\le m),\quad t\ge0,\quad
 q^\T u+\eta t<0.
\]
If \(t>0\), the point \(z=-u/t\) is feasible and \(q^\T z>\eta\).
If \(t=0\), the direction \(v=-u\) satisfies \(Bv\le0\), \(q^\T v>0\).
Choose a feasible base \(z_0\in Q\). Then \(z_0+s v\in Q\) for every
\(s\ge0\), and for large \(s\) its objective exceeds \(\eta\).
Both cases contradict validity. Hence
\[
 (q,\eta)=\sum_{i=1}^m\lambda_i(a_i,\beta_i)+r_0(0,1),
 \qquad \lambda\ge0,\quad r_0\ge0.
\]
The first coordinates give \(M\lambda=q\); the last coordinate gives
\(\beta^\T\lambda=\eta-r_0\le\eta\).
\end{proof}

The scalar \(r_0\) is an unused bound budget. It is helpful in this proof;
we need not introduce it as a variable in the eventual search.

Nonemptiness is indispensable. In variables \((x,y)\), the system
\(x\le0\), \(-x\le-1\) is empty and therefore implies \(y\le0\) vacuously.
Yet no combination of its normals can produce \((0,1)\).
There is a contradiction certificate, with weights \((1,1)\), but no bound
certificate with the requested objective normal. The correct first question
is which alternative holds, not whether a vacuous implication can be given
an arbitrary normal.

\section{Why a bound needs \texorpdfstring{\(k\)}{k} rows and a contradiction may need \texorpdfstring{\(k+1\)}{k+1}}
\label{sec:sparsity}

An immediate application of conic reduction in the lifted space would give
a \(k+1\)-generator bound, counting the vertical residual. Feasibility lets
us do better for bounds by reducing only the normal support while allowing
the combined right side to improve.

\begin{theorem}[Sparse bound certificates]\label{thm:sparse-bound}
Assume \(Q\ne\varnothing\) and \eqref{eq:reduced-certificate} has a solution.
There is a solution whose positive-support normal columns \(a_i\) are
linearly independent. In particular, at most \(k=\rank(B)\) inequalities
are needed.
\end{theorem}
\begin{proof}
Start with a certificate \(\lambda\). If its positive-support normals are
dependent, take a nonzero \(h\) on that support with \(Mh=0\).
We claim that its sign can be chosen so that
\begin{equation}\label{eq:orient}
 \beta^\T h\ge0\quad\text{and}\quad h_i>0
 \text{ for at least one }i.
\end{equation}
If \(\beta^\T h=0\), either sign with a positive entry works.
Otherwise choose its sign so that \(\beta^\T h>0\).
If all entries were then nonpositive, \(w=-h\ge0\) would satisfy
\(Mw=0\), \(\beta^\T w<0\). A feasible \(z_0\) would give
\(0=w^\T Bz_0\le w^\T\beta<0\), a contradiction.
Thus \eqref{eq:orient} holds.

Use the same minimum-ratio step \eqref{eq:conic-step}. It preserves
nonnegativity and normal matching, reduces support, and satisfies
\[
 \beta^\T(\lambda-t h)=\beta^\T\lambda-t\,\beta^\T h
       \le\beta^\T\lambda\le\eta.
\]
Iteration ends on independent normal columns, of which there are at most \(k\).
\end{proof}

Together with Theorem~\ref{thm:existence}, this proves the advertised bound
for every valid upper bound on a nonempty feasible set. The proof neither
assumes a vertex nor counts active rows. If \(q=0\), nonemptiness implies
\(\eta\ge0\) for validity, and \(\lambda=0\) suffices.

\begin{corollary}[Active rows certify a feasible optimizer]\label{cor:active}
Let \(z^\star\in Q\). It is an optimizer of \(q^\T z\) over \(Q\) if and
only if there is a nonnegative certificate supported on rows active at
\(z^\star\), with \(M\lambda=q\).
Such a certificate may be chosen with at most the rank of the active
normal columns.
\end{corollary}
\begin{proof}
If \(z^\star\) is optimal, Theorem~\ref{thm:existence} supplies a certificate
for \(\eta=q^\T z^\star\). Weak duality gives
\(\beta^\T\lambda\ge q^\T z^\star\), so equality holds.
The reduced gap identity
\[
 0=\sum_{i=1}^m\lambda_i(\beta_i-a_i^\T z^\star)
\]
forces every positively weighted row to be active.
On active rows, a dependence \(Mh=0\) also has
\(\beta^\T h=(z^\star)^\T Mh=0\), so conic reduction preserves the cost
exactly and achieves the active-rank bound.
Conversely, weights supported on active rows have
\(\beta^\T\lambda=\sum_{i=1}^m\lambda_i a_i^\T z^\star=q^\T z^\star\);
they certify a bound attained by the feasible point.
\end{proof}

This is the active-normal picture behind classical LP optimality
\cite[pp.~8--9]{bertsekas10}. It is not circular to search active rows at
a proposed point: first check feasibility, then ask whether the point has
a certificate at \emph{its own} objective value. Success proves optimality.
The point must not be called an optimizer before that proof is supplied.

\begin{theorem}[Sparse contradictions]\label{thm:sparse-contradiction}
If \(Q\) is empty, there is \(\lambda\ge0\) with
\[
 M\lambda=0,\qquad \beta^\T\lambda=-1,\qquad
 |\supp_+(\lambda)|\le \rho,
 \quad \rho=\rank[\,B\ \beta\,]\le k+1.
\]
The positive-support \emph{lifted} columns \((a_i,\beta_i)\) may be chosen
independent. In the infeasible case, \(\rho=k+1\).
\end{theorem}
\begin{proof}
Theorem~\ref{thm:farkas} supplies the representation
\((0,-1)=\sum_{i=1}^m\lambda_i(a_i,\beta_i)\).
Apply Lemma~\ref{lem:conic} to these lifted columns. Their span has dimension
\(\rho\le k+1\). Projection of that span onto its first \(d\) coordinates
has image of dimension \(k\). Its kernel lies on the one-dimensional vertical
axis and contains \((0,-1)\), so its kernel has dimension exactly one.
Rank-nullity gives \(\rho=k+1\).
\end{proof}

A nonzero contradiction necessarily has \emph{dependent} normal columns,
since their nonnegative combination is zero. Requiring normal independence
in contradiction search would rule out every nonzero certificate.
Independence belongs in the lifted space for this task.

\subsection{Sharpness and what the rank does not measure}

For the \(k\) constraints \(z_j\le1\), \(1\le j\le k\), certifying
\(\sum_{j=1}^k z_j\le k\) forces each multiplier to be \(1\).
Thus \(k\) rows really may be needed for a bound.

For infeasibility, take
\[
 z_j\le0\quad(1\le j\le k),\qquad
 -\sum_{j=1}^k z_j\le-1.
\]
Unit weights give \(0\le-1\). Omitting the last row permits \(z=0\);
omitting row \(j\) permits \(z_j=1\) with all other coordinates zero.
Every proper subsystem is feasible, so all \(k+1\) rows are necessary.
When \(k=0\), the final constraint is simply \(0\le-1\).

The rank \(k\) is not automatically the affine dimension of \(Q\).
With no explicit equalities, \(x\le0\), \(-x\le0\) define a singleton of
affine dimension zero, yet certifying \(x\le0\) needs one inequality.
Here \(k=1\). Turning a pair of inequalities into an implied equality with a
freely signed multiplier changes the accounting: the derivation from
original inequalities must still be retained.

\section{Finite exact discovery over the rationals}
\label{sec:search}

Suppose all original data, including the target bound, are rational.
Gaussian elimination produces rational \(x_0,N,B,\beta,q,\eta\).
The existence proofs above use real geometry, but their sparse certificates
can be found by rational linear algebra.

\begin{proposition}[Rational sparse witnesses]\label{prop:rational}
Whenever the hypotheses of the sparse-bound or sparse-contradiction theorem
hold for rational data, the asserted certificate can be chosen rational.
The recovered equality multipliers can also be chosen rational.
\end{proposition}
\begin{proof}
A sparse bound certificate on independent normal columns is the unique
solution of \(M_S\lambda_S=q\). Choose an invertible square row minor of
\(M_S\); its inverse, and hence this unique solution, are rational.
The remaining rows impose consistency conditions already satisfied by
the real solution. Apply the same argument to independent lifted columns
and target \((0,-1)\) for a contradiction. Finally,
\eqref{eq:recover} is a consistent rational linear system and has a rational
solution by Gaussian elimination. Empty supports use the empty rational
vector and require no minor.
\end{proof}

No approximation or search over a denominator grid is needed. The uniqueness
here is conditional on an independent support; it is not uniqueness among all
certificates.

\subsection{A deterministic procedure}

Keep each projected row's original index. The following procedure is a
mathematical algorithm using ideal exact rational arithmetic.

\begin{enumerate}[label=\textbf{\arabic*.},leftmargin=*]
\item \textbf{Normalize and eliminate.}
Convert all inequality directions and the objective sense.
Test \(Ex=f\). If inconsistent, return the equality-only contradiction.
Otherwise compute \(x_0,N,B,\beta,q,\eta\).

\item \textbf{Inspect constant rows.}
If \(a_i=0\) and \(\beta_i<0\), that row alone is a contradiction; normalize
its weight to \(-1/\beta_i\) and recover equality weights.
If \(a_i=0\) and \(\beta_i\ge0\), it can be omitted from discovery.

\item \textbf{For a specified bound, search normal supports.}
Assuming nonemptiness, enumerate subsets \(S\) in increasing cardinality,
then lexicographic original-row order, including \(S=\varnothing\), up to
size \(k\). Skip dependent columns. Solve \(M_S\lambda_S=q\) exactly and
accept precisely when a solution exists, \(\lambda_S\ge0\), and
\(\beta_S^\T\lambda_S\le\eta\).
For an exactly checked rational feasible candidate \(z^\star\), one may instead search only
its active rows with target \(q^\T z^\star\) and use their rank.

\item \textbf{For a contradiction, search lifted supports.}
Enumerate subsets up to size \(\rho=\rank[\,B\ \beta\,]\), in increasing
cardinality and then lexicographic original-row order, including the empty subset.
Skip dependent lifted columns, solve
\[
 \sum_{i\in S}\lambda_i(a_i,\beta_i)=(0,-1),
\]
and accept precisely when the unique solution exists and is nonnegative.
There is no extra vertical residual column in this contradiction search.

\item \textbf{Recover and verify.}
Pad unused inequality weights with zero, solve for \(\mu\), and check the
original coefficient identity, original right side, and all required signs.
For an optimum, separately check the candidate in every original row and
check that its objective equals the certified bound.
\end{enumerate}

If nonemptiness is not known, the contradiction search can be performed first.
Exhaustive failure of that search proves feasibility by
Theorem~\ref{thm:farkas}, although it need not produce a feasible point.
A practical calculation often has a feasible point available immediately.

\begin{theorem}[Completeness of the ideal finite search]\label{thm:complete}
For rational data, the exhaustive contradiction search finds a certificate
if and only if the system is infeasible. On a nonempty feasible set, the
exhaustive bound search finds a certificate if and only if the specified bound
is valid. At an exactly checked rational feasible candidate,
exhaustive active-row search finds a certificate if and only if the candidate
is optimal. All three searches terminate in ideal exact arithmetic.
\end{theorem}
\begin{proof}
Soundness follows from weighted addition and original-coordinate recovery.
Theorems~\ref{thm:existence}, \ref{thm:sparse-bound},
\ref{thm:sparse-contradiction}, Corollary~\ref{cor:active}, and
Proposition~\ref{prop:rational} supply a rational witness on one of the
enumerated independent supports whenever the corresponding property holds.
That support's unique exact solve will recover the witness.
There are finitely many subsets, and each finite rational elimination and
comparison terminates.
\end{proof}

Cardinality-first direct search, stopped at its first success, minimizes the
number of positive inequality weights among eligible certificates.
Indeed, a smaller certificate could be reduced to independent normal support
without increasing its size, and would have been tested earlier.
A zero coefficient in the first accepted support would likewise give a
smaller accepted subset. This statement does not minimize the attained
upper bound, the size of coefficients, or a multiplier norm.

\subsection{Complexity and safe shortcuts}

Before rank pruning, the bound search has at most
\begin{equation}\label{eq:count}
 \sum_{j=0}^{\min(k,m)}\binom mj
\end{equation}
supports; contradiction search has at most
\(\sum_{j=0}^{\min(k+1,m)}\binom mj\).
For fixed \(k\), these are polynomial in \(m\), of degrees at most \(k\)
and \(k+1\), respectively. When rank grows with input size, the subset
family can be exponential. Each independent support requires a rational
rank/consistency solve and sign and cost checks. These subset counts are
not bit-complexity bounds or runtime predictions: coefficient lengths and
the arithmetic used for elimination matter too.

For hand work, start with zero and singleton supports. Use active rows when
the candidate and purpose justify that restriction. Positive proportional
dominance is another safe reduction: if
\[
 a_i=t a_j,\quad t>0,\qquad \beta_i\ge t\beta_j,
\]
then row \(j\) implies row \(i\). In a weighted combination, replace weight
\(\lambda_i\) by adding \(t\lambda_i\) to row \(j\), and delete row \(i\)'s
weight. The normal is unchanged and the cost does not increase.
For a contradiction it remains negative and can be renormalized afterward.
Retain the original index of the row kept.

Arbitrary normal dependence does not justify deletion. The two rows
\(x\le10\) and \(2x\le2\) point in the same direction, but retaining the
first and discarding the second loses the useful bound \(x\le1\).
A heuristic may reorder supports; a claim of completeness still needs a
full fallback over the appropriate family.

\section{Normalization is a choice of representation}
\label{sec:normalization}

A certificate has several different notions of size. Its number of
nonzero inequality weights is one. The magnitudes of its multipliers,
their denominators, and the tightness of its bound are others. Normalization
can change some of these without changing the underlying halfspaces.

\subsection{Row scaling and nonuniqueness}

Replacing an inequality \((a_i,\beta_i)\) by
\((\kappa a_i,\kappa\beta_i)\), \(\kappa>0\), preserves its halfspace.
Dividing its multiplier by \(\kappa\) preserves its contribution to a
certificate and preserves positive support. An equality may be scaled by
any nonzero scalar with the reciprocal change in its multiplier.
Negative inequality scaling must reverse the inequality direction; zero
scaling can destroy information.

For objective \(x\) and rows \(x\le1,\ 2x\le2\), every
\[
 \lambda=(1-2t,t),\qquad 0\le t\le\tfrac12,
\]
certifies the optimum \(1\). Both endpoints have minimum support one.
Moreover, \(\|\lambda\|_1=1-t\), so minimizing that norm favors the
second, scaled row. Dividing the second row by \(2\) makes the rows identical,
and now every nonnegative pair summing to one has the same norm.
A norm-based preference therefore depends on units.

Even the standard contradiction normalization does not select unique weights.
For \(x\le0,\ x\le0,\ -x\le-1\), the family
\[
 \lambda=(t,1-t,1),\qquad 0\le t\le1,
\]
has zero normal and right side \(-1\). Its norm is always \(2\); the two
endpoints both have minimum support two.
A preferred row order or a coefficient-size objective can select a
representative, but soundness itself supplies no unique optimal
normalization.

A fixed nonzero objective prevents arbitrary rescaling of all bound
multipliers: multiplying them by \(\alpha\) changes their normal from
\(c\) to \(\alpha c\). Clearing denominators in a written identity is valid
only if its objective side is scaled as well. Contradictions have zero
normal, so they admit arbitrary positive rescaling before fixing right
side \(-1\).

\subsection{An exact right side can demand an unnecessary row}

For \(x\le0,\ -x\le1\), the target \(x\le1\) is valid.
If we insist that the weighted right side equal \(1\), rather than merely
be at most \(1\), we must solve
\[
 \lambda_1-\lambda_2=1,\qquad \lambda_2=1.
\]
The result is \((2,1)\), requiring two rows although \(k=1\).
The direct search instead accepts \((1,0)\): it proves the stronger bound
\(x\le0\), which certainly proves \(x\le1\).

One can impose exact equality in a lifted search by adding the residual
column \(e=(0,1)\). This is mathematically sound, but the resulting support
objective counts the residual as another generator.
It is not the same as minimizing original inequality support.

To see the difference even at \emph{first success}, take rows
\(x\le3,\ x\le1\), and target \(x\le2\).
The augmented columns, in order, are
\[
 v_1=(1,3),\qquad v_2=(1,1),\qquad e=(0,1).
\]
No singleton represents \((1,2)\). Lexicographic pair search first tries
\(\{1,2\}\), which succeeds with
\[
 \tfrac12v_1+\tfrac12v_2=(1,2).
\]
It has two positive inequality weights. The later pair \(\{2,e\}\) gives
\(v_2+e=(1,2)\), with only one positive inequality weight.
Direct normal search tests singleton row \(1\), rejects its cost \(3>2\),
and accepts singleton row \(2\), whose cost is \(1\).
Thus minimum lifted-generator support and minimum original-row support
can differ. No claim of universally fewer operations follows from this
example.

\section{Further calculations and diagnostic examples}
\label{sec:examples}

\subsection{Fractions arise naturally}

Maximize \(x\) under
\[
 2x+y\le\tfrac73,\qquad x-y=\tfrac12.
\]
Set \(y=t\), so \(x=\tfrac12+t\). The inequality becomes
\(3t\le\tfrac43\), and the objective is \(\tfrac12+t\).
Matching the projected objective coefficient \(1\) against the row
coefficient \(3\) gives multiplier \(1/3\).
The remaining original normal is
\[
 (1,0)-\tfrac13(2,1)=\tfrac13(1,-1),
\]
so the equality multiplier is also \(1/3\). Therefore
\begin{equation}\label{eq:fraction-example}
 \tfrac13(2x+y)+\tfrac13(x-y)=x
 \le\tfrac79+\tfrac16=\tfrac{17}{18}.
\end{equation}
The reduced endpoint \(t=4/9\) gives
\((x,y)=(17/18,4/9)\), which satisfies both rows and attains the bound.
This calculation discovers the weights by a one-column solve, not a guess
from a list of small integers.

\subsection{A zero normal can contain the whole contradiction}

Consider
\[
 x-y=\tfrac12,\qquad 2x-2y\le\tfrac34.
\]
Using \(x=\tfrac12+t,\ y=t\), the inequality becomes
\(0t\le-\tfrac14\). Thus \(k=0\), but a contradiction uses one inequality.
A weight of \(4\) makes the reduced right side \(-1\). The original
inequality normal then becomes \((8,-8)\), canceled by equality weight
\(-8\):
\[
 4(2x-2y)-8(x-y)=0,\qquad
 4\cdot\tfrac34-8\cdot\tfrac12=-1.
\]
This is the \(k+1\) phenomenon in its smallest dimension.
Deleting every zero projected row, or dividing every row by the norm of its
projected normal, would destroy this example.

\subsection{An optimum without any vertex}

Maximize \(x\) in variables \((x,y)\), subject only to \(x=1\).
Every \((1,y)\) is feasible and optimal. There is no extreme point, because
\[
 (1,y)=\tfrac12(1,y-1)+\tfrac12(1,y+1).
\]
The equality weight \(1\) and candidate \((1,0)\) certify value \(1\)
with zero inequality weights. In the reduced coordinates the objective
normal is \(q=0\). The feasible set is unbounded, but the objective is
constant. Neither boundedness of the feasible set nor a primal vertex is
needed for certificate soundness or the sparse-support conclusions.

More generally, on a nonempty feasible set the zero objective has value
zero at every feasible point. Zero multipliers prove every bound
\(0\le\delta\) with \(\delta\ge0\). A negative target cannot be valid there.

\subsection{Why an arbitrary solution of a dependent system is not enough}

Suppose the normal columns are \(-1,1\), and the target is \(1\).
Then normal matching requires
\[
 -\lambda_1+\lambda_2=1.
\]
If a row-reduction routine sets the free coordinate \(\lambda_2=0\),
it returns \((-1,0)\), which is inadmissible. But \((0,1)\) is nonnegative,
and so is \((t,t+1)\) for every \(t\ge0\).
One particular inadmissible solution is not a proof that all solutions are
inadmissible. Independent-support enumeration avoids this ambiguity:
the second singleton has the unique positive solution \(1\).

Sparsity does not bound denominator size. For a positive integer \(D\), the
single constraint \(Dx\le1\) and objective \(x\) force the certificate
weight \(1/D\). No fixed finite denominator grid works for all inputs.
Likewise, under
\[
 x+\tfrac1{1000}y\le1,\qquad -x\le0,
\]
the objective \(y\) has weights \((1000,1000)\), giving bound \(1000\),
attained at \((0,1000)\).
The determinant of the normal matrix is \(1/1000\), not zero.
Two rows can be enough while exact arithmetic is large or numerical
rank determination is delicate.

\section{What the published experiments and Lean results establish}
\label{sec:evidence}

This paper is an ordinary mathematical exposition based on the published
repository snapshot
\[
 \texttt{e4f337c3aa52ec4c3f2aab078d7c91b32ee0015f}.
\]
The source study, worked examples, and evidence records are cited in
\cite{study,toolkit,round}. The snapshot integrates several components;
the discussion here concerns the LP certificate and sparse-discovery work.

\subsection{The finite experiment}

The sparse-discovery benchmark compares direct normal search with augmented
target search on thirteen built-in rational fixtures. Both use increasing
support size followed by input-index order. The recorded report contains
\(44\) runs in total. Selected attempted-subset counts are:
\begin{center}
\begin{tabular}{@{}lcc@{}}
\toprule
Bound fixture & Direct first / full & Augmented first / full\\
\midrule
Sharp support, \(k=3\) & \(8/8\) & \(12/16\)\\
Duplicate and dependent normals & \(5/16\) & \(5/42\)\\
Rational equality reduction & \(4/5\) & \(4/16\)\\
Useful singleton last & \(8/64\) & \(8/163\)\\
Useful singleton first & \(2/64\) & \(2/163\)\\
\bottomrule
\end{tabular}
\end{center}
The first count stops at the first accepted representation; the full count
exhausts the specified family. Attempted subsets include the empty set and
dependent supports that are subsequently skipped. Counts of accepted
representations can include zero coefficients and repeats of the same
mathematical certificate. None of these is a count of distinct proofs
or a timing result.

The source-bound independent audit records \(5{,}040\) one-dimensional toy
inputs, each checked by two exhaustive strategies, for \(10{,}080\)
toy searches. It also records \(13\) zero-dimensional searches, two
first-success counterexample searches, six original-coordinate identities,
and two affine parametrizations. The audit used a determinant-minor and
Cramer-rule oracle, separate from the benchmark's elimination routines.
Fourteen ordinary regression tests and byte-identical reproduction of the
stored benchmark report are recorded. These numbers describe different
finite checks and should not be added into a single proof count.

The toy inputs have two or three rows in one dimension, normals in
\(\{-1,0,1\}\), budgets in \(\{0,1\}\), target coefficients in
\(\{-2,-1,0,1,2\}\), and target budgets in \(\{-1,0,1,2\}\).
Zero is feasible for every input. The Python checks use assertions and require
ordinary execution without the \texttt{-O} flag.
The toy family is deliberately bounded, and the fixture program supplies
its affine parametrizations. It is not an arbitrary-input equality
elimination pipeline, general LP status classifier, or exported proof bridge.
The table and audit counts above are \emph{reported results from the cited
snapshot}. This exposition extends neither the tested family nor the formal
Lean scope. No runtime comparison with simplex was performed.
They illustrate representation and ordering effects, not general performance
superiority.

\subsection{Soundness, discovery, and interface trust}

The separate \texttt{lp-certificates} component contains Lean soundness
theorems for typed certificate checks: an accepted certificate establishes
the stated result over all real feasible points. Its typed input layer also
proves normalization of original mixed row directions and objective sense.
That is a substantial but specific statement about checked witnesses.

The sparse-support theorems and ideal-search completeness proof in this
paper are ordinary proofs; they have not thereby become Lean theorems.
The existing executable finder remains deliberately incomplete.
A returned, exactly checked witness is useful; a bounded search returning
\emph{unknown} says no more than that it found no witness within its limits.

The external JSON parser and literal transcription into typed input remain
a trust boundary without an end-to-end refinement theorem.
Dependency caches, the compiler/kernel, and host execution retain their
documented trust roles. Recorded local replay uses the same Lean kernel;
it provides no validation by an independent kernel.
Reading saved receipts or inspecting their hashes
does not run the Lean kernel again. The publication also records limits of
dependency-cleanliness checks and saved-process metadata validation.
The detailed report is the appropriate place for those implementation
conditions~\cite{round,toolkit}.

In particular, the formal package does not claim a general completeness or
termination theorem for its finder, general strong duality, or formal
converse/existence results. Our ordinary proof of a finite exhaustive
procedure does not silently discharge those formal obligations.
Numerical approximations must still be reconstructed and checked exactly;
a small residual is not an identity.

\section{A practical conclusion}

For a hand-sized rational problem, the most effective first question is
often: \emph{which few constraints could add up to the objective?}
The answer becomes clearer after the equality directions have been removed.
A useful working order is:
\begin{enumerate}[label=\arabic*.,leftmargin=*]
\item Write all signs and bounds explicitly, then solve the equalities.
\item Look for an immediate constant-row contradiction.
\item Find or check a feasible candidate when possible; reduce the objective
      and identify its active rows.
\item Match the projected objective with one row, then two, and so on.
      Keep the right-side test as an inequality.
\item Recover signed equality weights and write the final identity using the
      original rows.
\item Check feasibility and attainment separately before declaring an optimum.
\end{enumerate}

The structural reason this can be economical is now precise:
a nonempty feasible set permits cost-aware reduction to \(k\) independent
projected normals, while a contradiction lives in a lifted space and may
need \(k+1\) rows. Neither argument requires artificial variables, splitting
free variables into positive parts, or adding one slack variable per row.
This is a principled way to seek short proofs. It does not remove the
combinatorial cost of discovery or replace general LP algorithms on large
instances.

\begin{thebibliography}{9}
\small

\bibitem{border}
KC Border.
\emph{Alternative Linear Inequalities}.
September 2013, revised 15 October 2020.
Author-hosted notes, California Institute of Technology.
Theorems 3--5, p.~3; Proposition 35, p.~17.
\href{https://kcborder.caltech.edu/Notes/Alternative.pdf}
{Author PDF}.

\bibitem{bertsekas3}
Dimitri P. Bertsekas.
\emph{Convex Analysis and Optimization}, MIT 6.253, Spring 2012,
Lecture 3. Conic Carath\'eodory theorem and proof, PDF pp.~8--9.
\href{https://ocw.mit.edu/courses/6-253-convex-analysis-and-optimization-spring-2012/3e2a100d7fa327637ed8e31771a22c18_MIT6_253S12_lec03.pdf}
{MIT OpenCourseWare PDF}.

\bibitem{bertsekas10}
Dimitri P. Bertsekas.
\emph{Convex Analysis and Optimization}, MIT 6.253, Spring 2012,
Lecture 10. Linear Farkas lemma, LP duality, and active normals,
PDF pp.~7--9.
\href{https://ocw.mit.edu/courses/6-253-convex-analysis-and-optimization-spring-2012/9f80ea2051c10cc0f8cf3839983c381a_MIT6_253S12_lec10.pdf}
{MIT OpenCourseWare PDF}.

\bibitem{boyd}
Stephen Boyd and Lieven Vandenberghe.
\emph{Convex Optimization}, lecture slides, revised with Parth Nobel.
Equality elimination, slide 4.35, PDF p.~122; also PDF p.~331.
\href{https://web.stanford.edu/~boyd/cvxbook/bv_cvxslides.pdf}
{Author-hosted slides}.
See also their \emph{Convex Optimization}, Cambridge University Press, 2004,
and the \href{https://web.stanford.edu/~boyd/cvxbook/}{book website}.

\bibitem{study}
\emph{Sparse discovery of linear-programming certificates}.
Codex for Math repository, \texttt{lp-sparse-discovery}, 13 October 2026.
Study, examples, source notes, benchmark, and independent reviews.
\href{https://github.com/Drsunflowerherbert/Codex-for-Math/tree/e4f337c3aa52ec4c3f2aab078d7c91b32ee0015f/lp-sparse-discovery}
{Exact source snapshot}.

\bibitem{toolkit}
\emph{Rational LP certificate toolkit}.
Codex for Math repository, \texttt{lp-certificates}.
Mathematical notes, soundness scope, and publication trust limits.
\href{https://github.com/Drsunflowerherbert/Codex-for-Math/tree/e4f337c3aa52ec4c3f2aab078d7c91b32ee0015f/lp-certificates}
{Exact source snapshot}.

\bibitem{round}
\emph{Mathematics research round 13 October 2026}.
Codex for Math repository, publication report.
\href{https://github.com/Drsunflowerherbert/Codex-for-Math/blob/e4f337c3aa52ec4c3f2aab078d7c91b32ee0015f/publication/ROUND-2026-10-13.md}
{Source-bound publication report}.

\end{thebibliography}
\end{document}


